

\section{Introduzione}

% \begin{frame}
%  \frametitle{Introduzione}
%  \vspace{0.3cm}
%     \begin{center}
%      
%    \large La piattaforma numerica consiste in un \textbf{ambiente computazionale semplificato e versatile}
%    che permette l'integrazione e l'accoppiamento dei codici computazionali di \textbf{\textcolor{bleudefrance}{neutronica}} e
%    di \textbf{\textcolor{mygreen}{termoidraulica}}.
%    
%     \end{center}
% 
% 
%  \begin{columns}
%   \begin{column}{0.48\textwidth}
%   \begin{center}
%   \begin{figure}
%    \centering
% %    \includegraphics[scale=0.2]{../../thesis/figure/result_tred/pacman}
%    \includegraphics[scale=0.2]{../../thesis/figure/result_tred/morenew/tempassurda}
%   \end{figure}
%    
%   \end{center}
% 
%   \end{column}
%   \begin{column}{0.48\textwidth}
%   \begin{figure}
%    \centering
% %       \includegraphics[scale=0.30]{../../thesis/figure/result_tred/pacman}
%    \includegraphics[scale=0.2]{../../thesis/figure/result_tred/new/tred7-pw.png}
%   \end{figure}
%    % % % %      \begin{center}
% % % % %        \includegraphics[scale=0.3]{../../thesis/figure/result_tred/new/tred-reactor}
% % % % %    \end{center}
%   \end{column}
% 
%  \end{columns}
% 
% \end{frame}





 %%%%%%%%%%%%%%%%%%%%%%%%%%%%%%%%%%%%%%%%%%%%%%%%%%%%%%%%%
\subsection{Reattore Nucleare}
 \begin{frame}
 \vspace{0.4cm}
 \frametitle{ALFRED (Advanced Lead Fast REactor Demonstrator)}
%  \framesubtitle{Advanced Lead Fast REactor Demonstrator}
\begin{center}
\large
 Progetto europeo di reattore nucleare (300 MW$_{th}$) che ha il ruolo di dimostrare la 
 \textbf{fattibilità della tecnologia dei reattori veloci raffreddati a piombo}, 
 i quali hanno caratteristiche di \underline{affidabilità} e di \underline{sostenibilità}.
\end{center}
 \begin{center}
 \begin{figure}
 \centering
 \includegraphics[scale=0.33]{../../thesis/figure/reactor/alfred.png}
 \end{figure}
 \end{center}
 
 \begin{center}
 \begin{itemize}
  \item Combustibile ossidi misti di uranio e di plutonio (MOX)
  \item Refrigerante Piombo Liquido
 \end{itemize}
 \end{center}
 
 \end{frame}
 
 %%%%%%%%%%%%%%%%%%%%%%%%%%%%%%%%%%%%%%%%%%%%%%%%%%%%%%%%%
 \begin{frame}
  \frametitle{Caratteristiche dei Reattori Veloci}
  \vspace{0.4cm}
  \begin{center}
  \begin{itemize}
  \item Prevalenza di reazioni di \textbf{fissione veloce}
  \item Ogni evento di fissione veloce libera \textbf{più neutroni} rispetto alle fissioni termiche
%   \item Tuttavia le reazioni di fissione alle alte energie sono \textbf{meno probabili} rispetto alle basse energie 
  \item \textbf{Maggiori assorbimenti} da parte degli isotopi fertili che trasmutano in isotopi fissili
  \begin{equation*}
 ^{238}_{92}U + ^{1}_{0}n \, \longrightarrow ^{239}_{92}U  \xrightarrow{\beta^-}
 \, ^{239}_{93}Np  \xrightarrow{\beta^-} \, ^{239}_{94}Pu
 \end{equation*}
  \item \textbf{Maggiore compattezza} dovuta all'assenza del moderatore di neutroni
 \end{itemize}
 \begin{block}{Sostenibilità}
 Capacità di \textbf{\textcolor{red}{autofertilizzazione}} del combustibile (Breeder Reactor), con la possibilità di generare più combustibile di quello che brucia
 \end{block}
 
 \end{center}
 \end{frame}

 
 %%%%%%%%%%%%%%%%%
% % %  \begin{frame}
% % %  
% % %  \frametitle{Reattori veloci raffreddati al piombo}
% % %   \begin{minipage}[c]{0.48\linewidth} 
% % %  \begin{center}
% % %  \begin{figure}
% % %  \includegraphics[scale=2.2]{../../thesis/figure/reactor/lfr.jpg}
% % %  \end{figure}
% % %  \end{center}
% % %  \end{minipage}\hfill
% % %  
% % %  \end{frame}
 %%%%%%%%%%%%%%%%%%%%%%%%%%%%%%%%%%%
 \begin{frame}
 \frametitle{Vantaggi e svantaggi del piombo liquido come refrigerante}
 \vspace{0.5cm}
 % \begin{minipage}[c]{0.48\textwidth} 
 \begin{block}{Vantaggi}
 \begin{itemize}
 \item[\checkmark] Mantiene uno spettro veloce dei neutroni
 \item[\checkmark] Ottime proprietà termodinamiche
 \item[\checkmark] Alta temperatura di ebollizione ($T_{ebollizione}=1739 ^\circ C$)
 \item[\checkmark] Bassa reattività con acqua o aria (inerte)
 \item[\checkmark] Schermaggio delle radiazioni dovute a decadimenti
 \end{itemize}
 \end{block}
 
 %  \end{minipage}\hfill
 
% \begin{minipage}[c]{0.48\textwidth} 
 \begin{block}{Svantaggi}
 \begin{itemize}
 \item[\textbf{$\times$}] Elevata tendenza ad essere corrosivo nei confronti dell'acciaio
 \item[\textbf{$\times$}] Opacità del piombo causa problematiche per l'ispezione e il monitoraggio
 \item[\textbf{$\times$}] Sistema ausiliario per mantenere il piombo liquido ($T_{fusione}=327.5 ^\circ C$)
 \end{itemize}
 \end{block}
 % \end{minipage}\hfill
 \end{frame}
 
 \begin{frame}
 \frametitle{Scopo della tesi}
 
%  \vspace{0.8cm}

  \begin{columns}
  \begin{column}{0.5\textwidth}
  
 \begin{block}{}
 \begin{itemize}
 \item Generare le \textbf{sezioni d'urto macroscopiche} mediate sulla composizione e sull'energia, partendo dalle librerie dei dati nucleari
 \item Modellare il reattore nucleare ALFRED
 \item Implementare il modello all'interno della piattaforma numerica
 \item Consentire l'accoppiamento tra i codici di neutronica e i codici di termoidraulica
 \item Simulare il comportamento multifisico
 \end{itemize}
 \end{block}
 
  \end{column}
  \begin{column}{0.44\textwidth}
  \vspace{0.8cm}
  \begin{figure}
   \centering
%       \includegraphics[scale=0.30]{../../thesis/figure/result_tred/pacman}
   \includegraphics[scale=0.28]{../../thesis/figure/result_tred/new/tred-reactor}
  \end{figure}
   % % % %      \begin{center}
% % % %        \includegraphics[scale=0.3]{../../thesis/figure/result_tred/new/tred-reactor}
% % % %    \end{center}
  \end{column}

 \end{columns}
% 
 
\end{frame} 
 %%%%%%%%%%%%%%%%%%%%%%%%%%%%%%%%%%%%%%%%%%%%%%%%%%%%%%%%%
